\section{Philosophy Around the World}

\paragraph{Hermann Keyserling}

Once a prominent philosopher, \textbf{Hermann Keyserling}is largely forgotten today. His philosophy was oriented toward freedom, action, elitism, reaction, and anti-rationalism. Apparently Evola met him in person at some point, but took a dislike to him. I wouldn't be surprised since Keyserling was a bit pompous. After all, he had achieved some worldly success; in those days, philosophers like Keyserling and Henri Bergson could actually get large audiences for their lectures ... not necessarily that the people were smarter, but probably because the philosophers were better. Nevertheless, I've recently learned that these frequent debates in our time about religion and God provide a prosperous income for the debaters.

Although I have never read any of his philosophical works, many years ago I found his two volume work \emph{The Travel Diary of a Philosopher} at a used bookstore. He had made a journey from Europe through the Near East, then India, China, Japan, across the Pacific on a steamer, then across the USA, finally returning to Europe. I recall feeling some envy that he could do such a thing. Keyserling was born to the landed gentry, so he had an independent source of income to support such a journey. He was drawn to philosophy after reading \emph{The Foundations of the 19\textsuperscript{th} Century}, by Houston Chamberlain. That would set him on the path of German Idealism, with the grand style of covering vast swaths of history and geography in his writings.

It is not easy to summarize the travel diary, since it is a collection of disjointed impressions without any sustained argument. Keyserling was actually quite sensitive to the best features of the cultures he encountered. In particular, he was interested in the philosophical and religious traditions of the countries he visited. However, his interest was not in the practices, dogmas, and metaphysics of those traditions, but rather in their effect on their adherents. We may say, in other words, that his interest was in the races of the soul of the various peoples.

\paragraph{The American Ideal}


Keyserling had mixed feelings after his 3000 mile journey across America. His comments are probably why American exceptionalists despise European intellectuals. For example, he writes:

\begin{quotationx}
Not freedom so much as arbitrariness predominates in America ---the arbitrariness, not of one, as in Asiatic despotism, but of every individual, which is not better.

The will of the people expresses itself on the whole as the rule of incompetence.

The emancipation of the spirit has led first of all to the rejection of all inherited wisdom---to immorality, positivism, nihilism---philosophies which are infinitely more foolish than those handed down from times of greater limitation.
\end{quotationx}


\paragraph{Erecting Barriers}


He points out that Americans desire freedom without understanding what freedom really is. Keyserling points out that Traditional order reflected a divine reality. The differences between men are real and cannot be legislated away. Overturning all traditional order can only lead to self-destruction. Perhaps this, and subsequent passages, gave Evola the idea for his metaphysical positivism.

\begin{quotationx}
The desire for freedom is awakened generally before recognition is ripe, before it is aware of what it means, how it is expressed, and this necessarily brings with it temporary coarsening and superficiality. The New World illustrates this state of things with terrifying clarity. The Americans have understood less than anyone else that, \emph{if the barriers which were imposed from without are to be removed, this is not for the purpose of dispensing with barriers altogether, but that they must give way to others which have been chosen freely}. They do not want it as yet to be true that the traditional orders among men, however conventional in detail, express realities; that differences in the age of the soul, in character, in talent, and even in inherited position, are something just as real as the differences between chemical elements, and that no God, as long as he remains in the sphere of nature, can act contrary to its laws; they want to be free without taking empirical reality into account.

The consequence is that life, instead of becoming more self-determined in a wider frame, loses it autonomous nature progressively. In the most modern democracy actions are determined mechanically to a degree which never existed under ancient tyranny: there, at any rate, something living took the decision, good or evil; here accident decides, the force of circumstances, the conjunction of events; here life is absolutely dependent on inorganic forces, as the ignorant chemist is dependent on the sweet will of his ingredients; if an explosive is compounded by his blind hands, he is blown up. But this experience had to be gone through.
\end{quotationx}


\paragraph{The Future of the World}


Yet, he feels that the world may become more like America, if only as the battleground for the triumph of the spirit:

\begin{quotex}
The spiritual part of man is strengthened in battle, and unfolds itself all the more completely and freely, the more resistance is overcome. Thus our present materialism is verily the guarantee for our future spirituality. The suitable body for this has already been prepared in America of today. Humanity of tomorrow will undoubtedly live in an external condition which will resemble most closely that of the United States. It will not recognize any rigid forms, and will allow absolute self-determination to everyone. It will, in rising above all nature, and in taking account only of what springs from the spirit, even realize the ideal of equality.
\end{quotex}


\paragraph{The Obscure Urges}


\begin{quotex}
Some day or other, democracy will have been overcome. Then, however, it will be revealed, to the surprise of many, that humanity was conscious once more, in following its obscure urges, of being on the right track. That lack of outer barriers which conditions tyranny and barbarism in America today will grant to a humanity of the highest inward culture the corresponding frame of life. Humanity will have acquired so much knowledge by then that its attitude to the soul will not differ from our attitude to nature. It will countenance psychic facts in exactly the same way as material ones; it will grant the man who inwardly is on a higher level the higher external position as a matter of course, without strife, well aware of the fact that it is just as senseless to decide a man's value by majority vote as to decide the question of the existence of Selenium.
\end{quotex}


This passage seems odd to me, since humanity has already acquired as much knowledge as is necessary. I am not so sanguine that the inner nature of men will be so visible, simply on the principle that the lower cannot understand the higher. Keyserling also claims that, at some point, the spirit will have become so powerful, that it can dispense with tradition. Perhaps, too, Evola accepts that, since the I, surpassing all his privations, freely chooses his representation of the world.

\paragraph{The Pragmatic Test}


Although Keyserling appreciates the wisdom of the East, believing it to be superior to the West, he sees them, the Indians in particular, as unsuitable for practical life. Europeans are more practical minded; rather than seeking to transcend the world of appearances, they instead strive to manifest their ideals in the physical world. That is the idealism of the West. He points out that the divergent religions in the West can nevertheless bring about material progress. In other words, the beliefs and dogmas seem to be secondary. He proposes the ``pragmatic test'', as the touchstone for the validity of an idea.

\begin{quotationx}
The spirit of Christianity is a spirit of practical work, for which reason it does not matter too much to what dogmatic concept it appears to be related at a given time.

The teachings of the Bhagavad Gita and Mahayana Buddhism are more profound than Christianity.

As a religion of practical action, it surpasses all others. Among Christian peoples alone the ideas of love, compassion, humanity have become objective forces, and this means that even the most imperfectly recognized metaphysical reality has been realized in appearance better through Christianity than through any other faith.

In the sphere of actual life \emph{the best version of an idea is the one which stands the pragmatic test best}---no matter to what extent it satisfies the mind. This is the meaning of that superiority of Christianity which history has proved, however much the one-sided intelligentsia may doubt it.

The practical superiority of Christianity is the expression of an absolute metaphysical advantage: it embodies, as no other religion does, the spirit of freedom. Man, conditioned by nature can show himself free only in two ways:

\begin{itemize}


\item By saying yes inwardly to all events

\item By taking the initiative in directing them.
\end{itemize}

Accordingly, two commandments sum up Christian ethics: that everyone should take his cross upon himself, and that everyone should fearlessly fight for the victory of good in a spirit of ready sacrifice. These two commandments really induce a life of freedom in everyone.

If the Indians, the profoundest of all thinkers, fail in practical life, this is due to the fact that they do not know how to impress their free being upon appearance.

\begin{itemize}


\item Instead of taking up their cross, they think of its insubstantiality.

\item Instead of letting the recognition of their essential unity with Brahma, who wishes to manifest himself more and more fully in this world, develop into action by displaying initiative everywhere in accordance with the Divine will, they merely watch how God helps Himself.
\end{itemize}
\end{quotationx}


\paragraph{The Sage and the Bodhisattva}


Keyserling concludes his diary with a meditation on the Sage and the Bodhisattva, concluding that the latter ``embodies the aim of human aspiration.''

\begin{quotationx}
The sage who, indifferent to the world, only strives after the realization of God. He is not beyond name and form as yet since it is \emph{he} who wants to see God.

The Bodhisattva has discarded this last bond. The Bodhisattva says yes to the most evil world, for he knows himself to be one with it. Rid of himself he feels his foundation in God, while his surface is intertwined with everything which exists.
\end{quotationx}



\flrightit{Posted on 2014-07-17 by Cologero}

